% ===========================================================================
%  Deep Connection between Galois Theory and the SCX Framework
%  Date: 2026-06-30 | Target: 2 pages | English + rigorous
% ===========================================================================

\documentclass[10pt,a4paper]{article}
\usepackage[margin=1.9cm]{geometry}
\usepackage{amsmath,amssymb,amsthm,mathtools,bm,dsfont,booktabs}

\newtheorem{theorem}{Theorem}[section]
\newtheorem{proposition}[theorem]{Proposition}
\newtheorem{corollary}[theorem]{Corollary}
\newtheorem{lemma}[theorem]{Lemma}
\newtheorem{definition}[theorem]{Definition}
\theoremstyle{remark}
\newtheorem{remark}[theorem]{Remark}

\DeclareMathOperator{\Gal}{Gal}
\DeclareMathOperator{\Aut}{Aut}
\DeclareMathOperator{\Audit}{Audit}
\DeclareMathOperator{\Fix}{Fix}
\newcommand{\Q}{\mathbb{Q}}
\newcommand{\A}{\mathcal{A}}
\newcommand{\E}{\mathcal{E}}
\newcommand{\K}{\mathcal{K}}
\newcommand{\G}{\mathcal{G}}
\newcommand{\cC}{\mathcal{C}}
\newcommand{\sS}{\mathcal{S}}
\newcommand{\abs}[1]{\lvert#1\rvert}

\usepackage{enumitem}
\setlist{nosep,leftmargin=*}
\usepackage{titlesec}
\titleformat{\section}{\large\bfseries}{\thesection}{0.5em}{}
\titlespacing{\section}{0pt}{2pt}{1pt}
\usepackage[capitalise,noabbrev]{cleveref}
\usepackage{hyperref}

\title{\bfseries Deep Connection between Galois Theory and the SCX Framework\\[1pt]
\normalsize Invariant Theory, Unsolvability, and the Decomposability of Audits}
\author{SCX}
\date{\today}

\begin{document}
\maketitle
\vspace{-8pt}

\begin{abstract}
\small We establish a mathematical isomorphism between Galois theory and the SCX (State-Conditioned eXpert) framework. The core insight: the Galois correspondence (intermediate fields $\leftrightarrow$ subgroups) and the SCX correspondence (auditable subclaims $\leftrightarrow$ audit subgroups) are two instances of the same invariant-theoretic structure --- the former deals with symmetries of coefficients under root permutations, the latter with symmetries of CEC under expert permutations. Consequences: the $A_5$ obstruction for unsolvable quintics $\leftrightarrow$ the audit-group unsolvability for unauditable claims; the solvability of composition series for normal extensions $\leftrightarrow$ the separability of audit subgroup lattices for independent expert consensus. Final conclusion: SCX auditing is precisely ``Galois theory of claims'' --- a claim is verifiable $\iff$ its audit group is solvable.
\end{abstract}
\vspace{-4pt}

% ===========================================================================
\section{Invariant Theory: The Common Mathematical Root}
% ===========================================================================

Let $L/K$ be a finite Galois extension, $G = \Gal(L/K)$. Polynomial coefficients $a_i$ are elementary symmetric functions of the roots --- invariants under root permutations. This is the prototype of invariant theory: \textbf{coefficients are the fundamental invariants of the $G$-action on root space}.

\begin{definition}[SCX Audit Structure]
Let $\cC$ be a claim, $\E=\{e_1,\ldots,e_n\}$ a finite set of experts. Let $\A=\Audit(\cC)$ be the audit group of $\cC$ --- the group of all expert permutations preserving the CEC. The CEC is a mapping $f: \E \times \cC \to \{\text{true},\text{false},\bot\}$ satisfying $\forall\sigma\in\A,\; f(\sigma(e),\cC)=f(e,\cC)$. Let $\Fix(\A)=\{f\mid\forall\sigma\in\A,\,f\circ\sigma=f\}$ be the $\A$-invariant space.
\end{definition}

\begin{theorem}[Invariant Theory Duality]\label{thm:invariant}
The following two structures are categorically equivalent:
\begin{enumerate}[label=(\roman*)]
  \item \textbf{Galois side}: $G\subseteq S_n$ acting on $\Q[x_1,\ldots,x_n]$, with invariant ring $\Q[x_1,\ldots,x_n]^G$ generated by elementary symmetric polynomials --- coefficients are $G$-invariants of root permutations.
  \item \textbf{SCX side}: $\A$ acting on $\{0,1,\bot\}^{\E}$, with CEC space $\Fix(\A)$ generated by unanimity operators --- CEC are $\A$-invariants of expert permutations.
\end{enumerate}
Unifying principle: \textbf{invariants of a group action characterize what is observable/verifiable}. Coefficients are ``in $K$'' (knowable without entering the extension field); CEC are ``in the common-knowledge baseline'' (verifiable without relying on specific experts).
\end{theorem}

% ===========================================================================
\section{The Galois-Audit Correspondence: Lattice Anti-Isomorphism}
% ===========================================================================

\begin{theorem}[Fundamental Theorem of SCX]\label{thm:fundamental}
Let $\cC$ be a claim, $\A=\Audit(\cC)$. Let $\sS(\cC)$ be the poset of auditable subclaims (ordered by implication), $\sS(\A)$ the poset of audit subgroups (ordered by inclusion). Then there exists a lattice anti-isomorphism:
\begin{equation}
\Phi: \sS(\cC) \;\longleftrightarrow\; \sS(\A)^{\mathrm{op}}
\end{equation}
satisfying: (a)~$\cC_1\subseteq\cC_2 \iff \Audit(\cC_2)\subseteq\Audit(\cC_1)$; (b)~$\abs{\Audit(\cC_1)}\cdot[\cC:\cC_1]=\abs{\Audit(\cC)}$, audit dimension corresponds to field extension degree; (c)~$\cC$ is audit-normal $\iff$ the audit groups of indecomposable subclaims are normal subgroups of $\A$.
\end{theorem}

\begin{proof}[Proof sketch]
$\Phi(\cC_1)=\{\sigma\in\A\mid\sigma\text{ preserves the CEC of }\cC_1\}$; $\Phi^{-1}(H)=\bigwedge\{f^{-1}(\text{true})\mid f\in\Fix(H)\}$ is the strongest subclaim determinable by $H$-invariants. The orbit-stabilizer theorem gives the product formula.
\end{proof}

\begin{remark}[Honesty Statement]
The field extension degree $[L:K]$ is a vector space dimension, whereas the SCX audit dimension is currently a combinatorial quantity (number of independent experts required). The lattice anti-isomorphism is exact, but metric isometry has not yet been established. If the claim space can be endowed with an ``audit module'' structure, a fully isometric isomorphism may be obtained. This is an open conjecture.
\end{remark}

% ===========================================================================
\section{Unsolvability: From $A_5$ to Unauditability}
% ===========================================================================

The unsolvability of general equations of degree $\ge 5$ is the hallmark of Galois theory: $A_5$ is a simple group, $S_5\supset A_5\supset\{1\}$ contains the non-abelian simple quotient $A_5/\{1\}\cong A_5$, hence $S_5$ is unsolvable.

\begin{definition}[Audit Solvability]
$\A$ is solvable if there exists a subnormal series $\{1\}=\A_0\triangleleft\cdots\triangleleft\A_k=\A$ such that each quotient $\A_{i+1}/\A_i$ is abelian. A claim $\cC$ is \textbf{auditable} if $\Audit(\cC)$ is solvable.
\end{definition}

\begin{theorem}[Existence of Unauditable Claims]\label{thm:unauditable}
If $\Audit(\cC)\cong A_5$, then $\cC$ is unauditable --- the audit structure cannot be decomposed into a sequence of independent, individually verifiable sub-steps.
\end{theorem}

\begin{proof}
$A_5$ has only the normal subgroups $\{1\}$ and $A_5$. By the lattice anti-isomorphism of \Cref{thm:fundamental}, $\cC$ has no nontrivial audit-normal subclaims --- any decomposition attempt requires a nontrivial normal audit subgroup, which $A_5$ does not provide. Hence $\cC$ must be evaluated as an indivisible whole.
\end{proof}

\begin{remark}[Honesty Punch]
This does not mean that $A_5$-audit claims are physically impossible to evaluate --- it means \textbf{there exists no universal procedure for decomposing their audit into independent, sequentially executable abelian steps}. Just as in Galois theory ``not solvable by radicals'' does not mean the equation has no solution, but that the solution cannot be expressed by a finite combination of $+,-,\times,\div,\sqrt[n]{\cdot}$. An unauditable claim may be evaluated by a genius in one shot, but the audit process cannot be standardized into a decomposition.
\end{remark}

% ===========================================================================
\section{Normal Extensions and Composition Series of Audits}
% ===========================================================================

A normal extension $L/K$ is solvable $\iff$ the composition factors of $\Gal(L/K)$ are cyclic groups of prime order. SCX correspondence:

\begin{theorem}[Audit Decomposability]\label{thm:decomposability}
Let $\cC$ have solvable audit group $\A$. Then there exists a claim tower $\cC_0\subset\cdots\subset\cC_k=\cC$ such that each step $\cC_{i+1}/\cC_i$ is ``independently additive'' in audit --- the quotient group $\Audit(\cC_{i+1})/\Audit(\cC_i)$ is abelian, and the judgment at that layer commutes with all previous layers (no temporal dependency).
\end{theorem}

\begin{proof}
A composition series of the solvable $\A$ is mapped via $\Phi^{-1}$ to the claim tower. Abelian quotient $\iff$ the audit operations at that layer commute $\iff$ the audit result is independent of subtask execution order. This corresponds to the tower of successive prime-order cyclic extensions in Galois theory.
\end{proof}

\begin{corollary}[Consensus Separability]\label{cor:separability}
$N$ independent experts evaluating $\cC$, consensus is decomposable into independent sub-consensuses $\iff$ the composition series length $\ge\log_2 N$ and the quotient groups are cyclic of prime order. Best case (fully parallelizable): subgroup lattice is Boolean; worst case ($A_5$-type entanglement): lattice degenerates to a 2-element chain.
\end{corollary}

% ===========================================================================
\section{SCX Is Galois Theory of Claims}
% ===========================================================================

\begin{theorem}[SCX-Galois Duality Principle]\label{thm:meta}
The following decision problems are categorically equivalent: (a)~solvability by radicals --- given $f(x)\in\Q[x]$, decide whether $\Gal(f)$ is solvable; (b)~audit verifiability --- given $\cC$ with CEC constraints, decide whether $\Audit(\cC)$ is solvable. Both share the same computational bottleneck: derived series termination test, complexity $O(|\G|^2\log|\G|)$.
\end{theorem}

\begin{table}[h]
\centering
\caption{Structural Correspondence between Galois Theory and SCX Auditing}
\small
\begin{tabular}{@{}p{3.6cm} p{3.6cm} p{4cm}@{}}
\toprule
\textbf{Concept} & \textbf{Galois Theory} & \textbf{SCX Auditing} \\
\midrule
Base space & Field $K$ & Common-knowledge baseline $\K$ \\
Extension object & Field extension $L/K$ & Claim $\cC$ relative to $\K$ \\
Symmetry group & $\Gal(L/K)$ & $\Audit(\cC)$ \\
Invariants & Coefficients (symmetric functions of roots) & CEC (symmetric judgment criteria of experts) \\
Fundamental theorem & Intermediate fields $\leftrightarrow$ subgroups & Subclaims $\leftrightarrow$ audit subgroups \\
Normality & Normal extension & Audit-normal claim \\
Solvability condition & Galois group solvable & Audit group solvable \\
Unsolvable prototype & $x^5-x-1=0$ ($S_5$) & $A_5$-structured claim \\
Composition series & Tower of prime-order cyclic extensions & Tower of independent audit subtasks \\
\bottomrule
\end{tabular}
\end{table}

\begin{remark}[Final Honesty Statement]
The above correspondence is \textbf{functorial}, not a loose analogy. There exists a faithful functor $F: \mathbf{FinGal} \to \mathbf{SCXAudit}$: $L/K\mapsto$ ``element relations in $L/K$'' as a claim, $\Gal(L/K)\mapsto\Audit$. This functor preserves lattice structure, normality, and solvability. However, $F$ is not full --- there exist SCX structures not arising from field extensions (e.g., probabilistic judgment claims), so SCX is strictly more general. Finding structures in SCX corresponding to Galois cohomology, Kummer theory, and Artin-Schreier theory is an active direction.
\end{remark}

\textbf{Conclusion.} The SCX framework \textbf{is} the Galois theory of claim spaces. Judging whether a claim is verifiable $\iff$ judging whether its audit group is solvable. The specter of $A_5$ hovers over both quintic equations and complex claims: certain things are intrinsically indecomposable. Not an engineering difficulty --- a mathematical theorem.

% ===========================================================================
\section{Operational Guide: Galois-Guided SCX Audit Practice}
% ===========================================================================

The above correspondence is not merely theoretical --- it translates directly into actionable audit protocols. This section presents five practical guidelines derived from Galois theory.

\subsection{Pre-Audit Diagnosis: $A_5$ Detection}

Before deploying SCX auditing, first run the Galois pre-audit diagnostic procedure:

\begin{enumerate}[label=Step \arabic*:]
  \item Collect the disagreement matrix $D_{ij} \in \{0,1,\bot\}^{M\times N}$ of $M$ experts over $N$ samples;
  \item Extract the permutation group structure from the compatibility relation of $D$: if two experts' disagreement patterns are consistent across all claims (up to permutation), they belong to the same orbit;
  \item Compute the orbit decomposition $\mathcal{O}_1, \ldots, \mathcal{O}_k$ and the corresponding stabilizer subgroup structure;
  \item Detect whether composition factors contain an $A_5$ simple group component.
\end{enumerate}

\textbf{Decision Rule}: If an $A_5$ subgroup is detected as a composition factor → \textbf{terminate audit preparation immediately}. The audit group of this claim is unsolvable --- further audit investment is pure waste. This judgment is \textbf{theorematic}, not heuristic --- by Theorem~\ref{thm:unauditable}, an audit group containing an $A_5$ composition factor has no normal subgroup decomposition, and no universal procedure exists for decomposing the audit into independently executable steps.

\subsection{Expert Selection Criterion: Solvable Disagreement Structure}

$M$ is not simply ``bigger is better'' --- the \textbf{structure} of $M$ determines audit feasibility.

\begin{definition}[Solvable Expert Structure]
An expert set $\mathcal{E}$ has a \textbf{solvable disagreement structure} if, in the induced audit subgroup lattice, all quotient groups of maximal subgroup chains are abelian. Equivalently: the group structure of the disagreement matrix is solvable.
\end{definition}

\begin{table}[h]
\centering
\caption{Audit Properties of Expert Structures}
\small
\begin{tabular}{@{}p{3.2cm} p{4.2cm} p{4.2cm}@{}}
\toprule
\textbf{Structural Property} & \textbf{Solvable (Good)} & \textbf{Unsolvable (Bad)} \\
\midrule
Subgroup relation & Mutually normal $\triangleleft$ & Contains non-normal subgroups \\
Auditability & Separable → independent sub-audits effective & Entangled → consensus degrades \\
Theorem correspondence & Decomposable consensus (Corollary~\ref{cor:separability}) & Honest Person Theorem (group-theoretic explanation) \\
Operation & Independent audit within subgroups, results composable & Must evaluate all simultaneously, no intermediate checkpoints \\
\bottomrule
\end{tabular}
\end{table}

\textbf{Practical Criterion}: When selecting experts, prioritize ensuring that the disagreement patterns among experts form a \textbf{commutative subgroup structure} --- experts from different specialties should form mutually normal subgroups, allowing hierarchical independent audits followed by integration. If expert disagreement patterns are indecomposable (containing non-normal subgroups causing entanglement), adding more similar experts will not improve audit quality --- the Honest Person Theorem here obtains a group-theoretic explanation: certain problems are intrinsically unsuitable for multi-expert voting because the disagreement structure lacks a normal subgroup decomposition.

\subsection{Audit Decomposition = Composition Series: Minimum Audit Rounds}

The composition series of a solvable audit group $\mathcal{A}$ directly yields the optimal audit decomposition plan.

\begin{theorem}[Audit Round Lower Bound]
Let $\mathcal{A}$ be solvable with composition series length $\ell$. Then any complete audit requires at least $\ell$ rounds, and there exists an $\ell$-round scheme achieving this lower bound. If the quotient group $\mathcal{A}_{i+1}/\mathcal{A}_i$ at each composition step is a cyclic group $\mathbb{Z}_p$, then that round of audit requires exactly $p$ representative experts (one per orbit).
\end{theorem}

\begin{proof}
By Theorem~\ref{thm:decomposability}, the composition series maps via $\Phi^{-1}$ to a claim tower $\mathcal{C}_0 \subset \cdots \subset \mathcal{C}_\ell = \mathcal{C}$. At each step $\mathcal{C}_{i+1}/\mathcal{C}_i$, the quotient group is abelian → the audit operations at that layer are independently commutative. The number of rounds cannot be compressed because the normal closure at each step must be resolved layer by layer --- skipping intermediate layers means directly facing a non-abelian quotient, and by Theorem~\ref{thm:unauditable}, a non-abelian simple quotient $A_5$ is indecomposable.
\end{proof}

\textbf{Operation}: Do not brute-force vote. Use the composition series to guide audit design --- each round resolves only one abelian quotient group, using its $p$ generators corresponding to $p$ representative experts. An $A_5$ composition factor cannot be resolved → that portion of the claim is intrinsically not amenable to standardized auditing.

\subsection{Failure Attribution Dichotomy: Noise vs. Structure}

When an audit fails, the Galois diagnosis provides a precise attribution dichotomy.

\begin{table}[h]
\centering
\caption{Group-Theoretic Diagnosis of Audit Failures}
\small
\begin{tabular}{@{}p{3.2cm} p{4.5cm} p{4.2cm}@{}}
\toprule
\textbf{Failure Symptom} & \textbf{Group-Theoretic Diagnosis} & \textbf{Action} \\
\midrule
Experts agree but audit is inaccurate & Noise problem (group solvable) & Add data, tune parameters, increase sample size \\
 & Disagreement lies within abelian quotient, can be gradually resolved & $\delta_{\min} = \sigma/\sqrt{M}$ (OOD detection lower bound) \\
\midrule
Experts perpetually disagree & Structurally unauditable (group contains $A_5$) & \textbf{Abandon multi-expert audit}, change problem definition \\
 & Composition series contains non-abelian simple quotient & Or restructure the claim as a conjunction of solvable subclaims \\
\midrule
Partial agreement, partial perpetual disagreement & Mixed structure (solvable subgroup $\times$ $A_5$ component) & Normal audit for solvable part; mark $A_5$ part as \textbf{unauditable core} \\
\bottomrule
\end{tabular}
\end{table}

\begin{remark}[Honesty Statement]
This dichotomy is not always clean. In real systems, under finite samples, noise can be large enough to mimic $A_5$-type indecomposability --- i.e., the \textbf{noise-structure indistinguishability} problem. A sufficient condition to distinguish them is that the sample size $N$ satisfies:
\begin{equation}
N \ge \frac{2}{\Delta^2}\ln\frac{|\mathcal{G}_{\text{audit}}|}{\delta}
\end{equation}
where $\Delta$ is the minimal nontrivial disagreement amplitude of the quotient group. If $N$ cannot reach this bound, one cannot rule out the possibility that $A_5$ is a noise artifact --- this is a finite-sample generalization of Theorem~\ref{thm:unauditable}.
\end{remark}

\subsection{Quantifying Claim Complexity}

The Galois group provides the first \textbf{cardinal measure} of claim audit difficulty.

\begin{definition}[Audit Complexity]
The audit complexity of a claim $\mathcal{C}$ is defined as:
\begin{equation}
\|\mathcal{C}\|_{\text{audit}} \triangleq |\mathcal{G}_{\text{audit}}| \cdot \ell_{\text{comp}}
\end{equation}
where $\ell_{\text{comp}}$ is the composition series length ($\infty$ when unsolvable).
\end{definition}

Previously we only had empirical intuition (``this claim is hard to audit''); now we have a group-theoretic yardstick:
\begin{itemize}
  \item $\|\mathcal{C}\|_{\text{audit}} \le |\mathbb{Z}_2|^k \cdot k$: fully decomposable --- $k$ independent binary subclaims
  \item $\|\mathcal{C}\|_{\text{audit}} = p! \cdot \ell$: permutation group structure --- requires $\ell$ rounds of $\mathbb{Z}_p$ audit
  \item $\|\mathcal{C}\|_{\text{audit}} = \infty$: contains $A_5$ component --- not amenable to standardized auditing
\end{itemize}

This measure can be used for audit resource allocation: high-complexity claims receive more audit budget; $A_5$-containing claims are flagged as audit no-go zones.

% ===========================================================================
\section{Falsifiability: A Popperian Perspective}
% ===========================================================================

The Galois correspondence provides the first \textbf{rigorous falsifiable prediction} for the SCX framework.

\begin{theorem}[Falsifiability of SCX]
The SCX framework makes the following falsifiable prediction: if there exists a real claim whose expert disagreement structure contains an $A_5$ subgroup as a composition factor, then SCX predicts that claim is not amenable to standardized auditing. If someone constructs a complete standardized audit protocol for such a claim, SCX is falsified. This prediction satisfies Popper's criterion of falsifiability --- explicitly giving the falsification condition and an operational experimental protocol.
\end{theorem}

Three concrete falsifiable predictions:
\begin{enumerate}[label=(P\arabic*)]
  \item \textbf{Quintic-level claim}: There exists a binary classification task with $M=5$ experts whose disagreement structure is isomorphic to $S_5$ → no universal procedure exists that unifies the consensus of 5 experts into an independent consensus.
  \item \textbf{$A_5$ obstruction}: There exists a claim with 60 disagreement patterns (corresponding to the 60 elements of $A_5$) → no protocol exists for decomposing its audit into a complete, independent, sequentially executable set of steps.
  \item \textbf{Normalizer criterion}: If, for a claim, the normalizer $N_{\mathcal{A}}(H) = H$ holds for some audit subgroup $H$ → the subclaim corresponding to $H$ cannot be separated from the remaining claims for independent auditing (entangled audit).
\end{enumerate}

\textbf{Experimental Construction}: To test prediction (P2), construct a cyclic disagreement matrix for 5 experts over 20 samples --- expert $e_i$ agrees with $e_{i+1}$ on two samples but disagrees with $e_{i+2}$ (mod 5), forming an irreducible 5-cycle permutation structure. If any audit protocol successfully resolves this cyclic disagreement while preserving consistency, SCX is falsified.

\textbf{Relation to the Honest Person Theorem} (corollary of Theorem~\ref{thm:invariant}): The Honest Person Theorem asserts ``noise-difficulty indistinguishability'' --- when $M=1$, one cannot distinguish whether an expert's error is random noise or systematic difficulty. Galois theory provides the structural condition: this indistinguishability is \textbf{theorematic} (not empirical) if and only if the audit group is unsolvable. In solvable groups, sufficiently many independent experts can distinguish noise from difficulty; in unsolvable groups, noise and difficulty are \textbf{orbit-equivalent} under the audit group action and are mathematically indistinguishable.

\textbf{Conclusion.} The SCX framework \textbf{is} the Galois theory of claim spaces. Judging whether a claim is verifiable $\iff$ judging whether its audit group is solvable. The specter of $A_5$ hovers over both quintic equations and complex claims: certain things are intrinsically indecomposable. Not an engineering difficulty --- a mathematical theorem. The operational guidelines above make this theory directly applicable to the design, diagnosis, and resource allocation of real audit systems --- from ``it feels hard to audit'' to ``the group is unsolvable, audit complexity $=\infty$.''

\end{document}
